\scanpage{96}
\mspage{89}
% archive sheet number: 95
\deleted{\textcircled{$\unsure{e'|e}$} $\underline M$ satisfait $L_\pi$, \deleted{\unsure{et}} le foncteur d'inclusion $\underline M\longrightarrow E$ est $\pi$-accessible $^\ast$[et admet \deleted{\textcircled{$\unsure{d|e|c}$}} une sous-catégorie strictement génératrice contenue dans $\underline M^{\pi}$].
% slanted
\margin{Cette condition \unsure{semble plus} \unsure{impliquée|compliquée}, \unsure{simplement|sûrement}.}}

Je présume que ces conditions sont équivalentes.

$b_\pi\Rightarrow a_\pi$. On suppose b) vérifié, on veut prouver a. D'abord que \deleted{tout} $X\in M$ est de la forme dite $\varinjlim_I X_i$, avec les $\unsure{X_i|\alpha_i}\in M^{\pi}$ et $I$ grand devant $\pi$. Par prop. 1', on sait qu'on a une présentation parfaitement adaptée pour $\pi$, $X=\unsure{\underline{\lim}|\varinjlim}_{i\in I}X_i$. Par b), \inserted{``il faut''} la partie $I_M$ de $I$ est cofinale à $I$, donc comme $I$ est grand devant $\pi$, $I_M$ aussi, et on aura (posant $J=I_M$)
\[
X=\varinjlim_{i\in I_M}X_i,\qquad\begin{aligned}
X_i&\in M^{\pi},\\
I_M&\text{ grand devant }\pi,
\end{aligned}
\]
OK. Inversement, prouvons que si $X$ se met sous la forme dite, \inserted{dans a),} il est dans $M$. Cela va résulter du

\uline{Lemme 1}: Si \deleted{\unsure{les}} la condition \deleted{\unsure{a)}} b) \deleted{\unsure{(on admet $a_\pi$, qui est impliqué par b))}} est satisfaite, $M$ est stable par limites inductives grandes devant \deleted{\unsure{$\pi_0$|$<$}} $\pi$.
% slanted
\margin{\deleted{\unsure{On peut} \unsure{prendre} \unsure{filtrant}}}

\deleted{\unsure{Pour le moment, on a}}

Il faut prouver que si $\unsure{\deleted{X}X_i|X_i|X}$
